\section{Proof of the neutral deformation theorem}\label{d:proof}

We now prove Theorem~\ref{d:deformation}.  Throughout this section the
data have already been strictly prepared as in that theorem, and we
write $K=k$ and $J_0=J_{\rm con}$.  Unless indicated otherwise, all
derivatives and contractions use $g$.  Several letters are local to
this proof: $h$ will be a rescaled graph height, $L$ a profile, $a$ a
gradient length, and $\chi$ an ellipticity coefficient.  They are
distinct from the meridional metric, crossing lower bound, model
rotation parameter and electromagnetic potential used earlier.

\paragraph{Order of choices and limits.}
\begin{enumerate}
\item Fix the prepared data, then the reduced domain, collars and cutoff shapes.
\item Fix a sufficiently small $\eps>0$, choose a sufficiently large $n$,
      and form its profiles.
\item At this fixed $n$, solve on sufficiently large truncations $\Omega_R$
      and let $R\to\infty$.
\item Pass to $n\to\infty$ in the final energy estimate; no limit of the
      deformation metrics is required. Section~\ref{d:reduction} subsequently
      sends $\eps\downarrow0$ and removes the preparation.
\end{enumerate}
Here $R$ is the outer truncation radius, distinct from the model radius in
 the comparison. Polynomial bounds $\Pi_n$ are kept separate from constants
 that may depend arbitrarily on fixed $(n,\eps)$.

\subsection{Full trapped regions and stable collars}\label{d:domains}

We use the following precise form of the barrier and trapped-region
theorems~\cite[Theorems 3.1, 4.1 and 7.3]{AnderssonMetzger2009}.
In a connected compact smooth three-dimensional data region, let the expansion
for a prescribed smooth tensor $Q$ be $H+\tr_\Sigma Q$, with the
normal out of a bounding inner region.  A nonempty strictly positive
outer barrier, together with a strictly negative required inner
barrier, gives a smooth stable enclosing marginally outer trapped
surface.  The designated outer barrier may be disconnected.  One
may also take no required inner barrier: in that case the trapped
region is allowed to be empty.  If nonempty, the full trapped region,
defined as the union of the inner sides of all smooth bounding
surfaces with nonpositive expansion and containing any required
inner barrier, has a smooth stable outer frontier; its closure is
itself a maximal such region.  Bounded complementary components
meeting no designated outer face may be filled.  For a disconnected
region apply the result separately to each component meeting a
designated outer face: the strict existence statement requires a
nonempty inner barrier in that component, whereas the trapped-region
statement allows an empty inner barrier and an empty trapped region.
No energy condition
on $Q$ is required.  All applications below are on compact
truncations with these barrier signs; they require no extension of
the original constraints across the inner boundary.

For $\max_{\partial\Omega}\theta_+<c<0$, let $\mathcal B_c$ be the
closed full black region for
$H+\tr_\Sigma K\le c$, with the original boundary prescribed as an
inner barrier.  This is the preceding construction for
$Q=K-(c/2)g$, because a two-dimensional tangential trace of
$(c/2)g$ is $c$.  These regions are nonempty and contain an inner
collar.  They lie in a fixed compact radial range: at the largest
radius of a bounding surface, comparison with a large coordinate
sphere, on which $K=0$ and $H>0$, rules out a negative expansion.
This also makes the full regions independent of a sufficiently
distant truncation.

Choose a threshold $c_b<0$.  In the exterior of its full black region,
form the closed full white regions $\mathcal W_c$ for
$H-\tr_\Sigma K\le c<0$.  There is no required inner face, and every
black boundary is a designated outer face.  Its reversed normal has
white expansion $-c_b>0$ before the shift.  Together with the distant
sphere, these are strict outer barriers for the tensor
$-K-(c/2)g$.  The white region may be empty; otherwise its smooth
stable frontier is disjoint from the black region.
The black construction takes place in the connected exterior
truncation with both barriers.  After bounded pockets are filled,
the white construction takes place in its connected complement
toward infinity, with all black faces and the distant sphere
designated as outer faces and with no required inner face.  Thus it
uses the trapped-region statement, not the strict existence
statement with a missing inner barrier.

Choose $c_b$, and then $c_w$, arbitrarily close to zero at points of
Hausdorff right continuity of their respective full families.  To
justify the choice, take the distance to each closed region, capped
by a number larger than the diameter of the common compact core.
At each point of a fixed countable dense set this is a monotone
function of the threshold and has at most countably many
discontinuities.  Away from the union of these exceptional sets,
the common Lipschitz bound upgrades pointwise right continuity to
uniform right continuity.  This is equivalent to Hausdorff right
continuity for the nonempty regions.  For the empty region it means
that regions at sufficiently close larger thresholds remain empty.
Filling bounded pockets makes the black complement connected toward
infinity for the subsequent white construction.  Invariance follows
because the families are full maximal regions, so the circle action
maps each family to itself.

Let $\Omega_{\rm red}$ be the closure of the component toward infinity
outside both selected regions.  Its smooth compact boundary
$B=B_b\cup B_w$ consists of whole frontier components and is nonempty,
since it separates the entire original obstacle from infinity.  A
color may be absent.  With $\nu$ pointing into $\Omega_{\rm red}$ and
$P_B=\tr_BK$, the boundary identities are
\begin{equation}\label{d:threshold-boundary}
 H+P_B=c_b\quad(B_b),\qquad H-P_B=c_w\quad(B_w).
\end{equation}

Each face has a smooth outward collar of leaves, with parameter
$s\ge0$, $|\nabla s|>0$, whose expansion in its own color is at least
the corresponding threshold.  Here is the treatment of possible
zero stability eigenvalues.  For the shifted tensor, the principal
stability eigenvalue is nonnegative.  If positive, its positive
eigenfunction gives an outward graph speed with positive expansion
derivative.  If it is zero, the augmented linear map
\[
 (v,d_0)\longmapsto
 \left(L_{\rm MOTS}v-d_0,\ \int_Bv\,\dd A\right)
\]
is an isomorphism on each component.  Indeed the positive principal
eigenfunction spans the kernel, the positive adjoint eigenfunction
spans the cokernel, the constant $d_0$ removes the latter, and the
integral condition removes the former.  The implicit function
theorem therefore gives constant-expansion graphs at prescribed
small positive $\int_Bv$.  Their initial speed is positive.  Their
expansion cannot be at most the threshold, since that would enlarge
the full maximal region.  These graphs give the desired leaves.
Invariant eigenfunctions and uniqueness in the augmented problem
make the collars invariant.  Shorten them so that all collars are
disjoint, and let $\nu_s=\nabla s/|\nabla s|$ be their outward normals.

Choose a smooth compactly supported invariant offset $C(x)$, constants
$C_w\le C\le C_b$, and a smooth invariant weight $\rho>0$, equal on
the far end to a sufficiently small multiple of $r^{-4}$, such that
\begin{gather}\label{d:assigned-heights}
 b_-=c_b-C_b<0,\qquad b_+=-c_w-C_w>0,\\
 |(h+C)\tr K|+|\nabla C|+\rho\le\frac{\mathfrak m}{2}
 \quad\text{whenever }b_-\le h\le b_+.
 \label{d:offset-margin}
\end{gather}
Near a black face require $C=C_b-k_Bs$, and near a white face require
$C=C_w+k_Bs$, with $k_B>0$.  These choices are possible in the
following order.  First take the thresholds sufficiently close to
zero, using compact support of $K$ and the positive compact margin.
After fixing the collars, take sufficiently small opposite offsets,
with the stated linear profiles cut off to zero.  Finally take
$\rho$ small on the compact part and use the end lower bound for
$\mathfrak m$ to choose its $r^{-4}$ tail.  For a missing color its
constant is only an upper or lower bound.

\begin{lemma}[Enclosure from exterior supports]\label{d:weak-support}
Work in a connected compact barrier truncation as above.  Let $D$ be
a nonempty compact closed set, disjoint from the designated outer
faces and containing a collar of any required inner face.  Suppose
that every smooth exterior support at an interior frontier point of
$D$ has expansion at most $c$ for a smooth tensor $Q$.  An exterior
support means a smooth regular level through the point that contains
$D$ locally on its inner side.  Then, for every $c'>c$ for which the
outer barriers remain strict, $D$ is contained in a smooth bounding
region with outward expansion at most $c'$.
\end{lemma}
\begin{proof}
First show that every exterior support of a small parallel set
$D_z=\{x:\dist(x,D)\le z\}$ has expansion at most $c+Cz$, with $C$
depending on the compact geometry and $Q$, but not on its curvature.
Let such a support touch at $x$, and take a shortest segment of length
$z$ from $y\in D$ to $x$.  Choose a local unit vector field $V$ agreeing
at $y$ with its initial tangent, and set $F(y')=\exp_{y'}(zV(y'))$.
Prescribe the first derivative of $V$ so that $\dd F_y$ is exactly
parallel transport along the segment.  For each initial vector this
is the Jacobi boundary-value problem with that vector and its parallel
translate as the two endpoint values.  Rescaling the Jacobi equation
shows that the required initial derivative is $O(z)$ times the
initial norm.  Its component along the segment is zero, since the
parallel component of the Jacobi field is constant.  Thus it is a
permissible first jet of a unit field.  More explicitly, if
$B_i=\nabla_iV$, the unit-length constraint on its symmetric second
jet is
$\langle\operatorname{sym}\nabla^2_{ij}V,V\rangle
=-\langle B_i,B_j\rangle$; its tangential symmetric part can be
chosen bounded, and the antisymmetric part is fixed by the bounded
curvature commutator.  In a smooth orthonormal frame, normalizing a
vector-valued polynomial with the prescribed value and first jet
realizes precisely this compatibility.  Thus $V$ has bounded second
derivatives independently of the support.  Since $F$ is the identity
at $z=0$, smoothness of the exponential map now gives
$|\nabla\dd F|=O(z)$.

If $\phi$ defines the support at $x$, then $\phi\circ F$ supports $D$
at $y$: every $y'\in D$ maps a distance at most $z$ from $D$.  The
tangent ball at $x$ aligns the support normal with the final segment
tangent.  Since $\dd F_y$ is exactly parallel transport, comparison
of the two tangential Hessian traces after gradient normalization
has error only $O(z)$, coming from $\nabla\dd F$; there is no
uncontrolled multiple of $\nabla^2\phi$.  Normal transport and
variation of $Q$ contribute another $O(z)$.  This proves the claimed
parallel-set bound.

Choose $0<a<b$ so small that $c+Cb<c'$.  Smooth approximation of the
distance followed by a regular-value choice gives a smooth enclosure
containing $D_a$ and contained in $D_b$; fill any pockets that meet no
designated outer face.  Let $\eta\ge0$ be one on its boundary and
supported where $\dist(\cdot,D)<b$.  Replacing $Q$ temporarily by
\[
 Q-\frac12(c'+A_0\eta)g
\]
makes this enclosure a strict inner barrier when $A_0$ is sufficiently
large.  The designated outer barriers stay positive.  Apply the
strict barrier theorem to each complementary component meeting a
designated outer face; every such component has both barriers, by
connectedness and the strict enclosure.  The resulting marginal
surfaces together bound a region containing $D_a$.

If the minimum distance $z$ of that boundary from $D$ were at most
$b$, its boundary at a minimizing point would be an exterior support
of $D_z$.  Indeed a path from $D$ of shorter length cannot cross that
boundary before reaching it.  Its expansion for the original $Q$
is $c'+A_0\eta\ge c'$, contradicting the strict parallel-set bound.
Thus the boundary lies beyond $D_b$, where $\eta=0$, and has original
expansion $c'$.  This proves the enclosure assertion without any
regularity hypothesis on $D$ or any support-curvature bound.
\end{proof}

\subsection{The equations and their geometric identities}\label{d:equations}

Fix $0<\eps<1$ sufficiently small and an integer
$n>\max\{4,2/\eps\}$.  Choose a smooth cutoff $\vartheta$ equal to
one on $(-\infty,0]$, zero on $[1,\infty)$, and between zero and one.
Define
\begin{equation}\label{d:profiles}
 p(t)=n\vartheta(nt),\qquad
 l(t)=\exp\left(\int_0^t p(s)\,\dd s\right),\qquad
 L(t)=l(t)e^{2t},\qquad
 \ell=e^{-n\eps},\qquad \tau_0=\ell^{3/2}.
\end{equation}
For a graph function $f$, put
\begin{gather}\label{d:variables}
 \begin{gathered}
 h=\tau_0f,\qquad \sigma=|\nabla f|,\qquad D=1+l^2\sigma^2,\qquad
 d=D^{-1},\\
 w=\frac{l\nabla f}{\sqrt D},\qquad
 a=|w|,\qquad v=a^2=1-d,\qquad u=L\sqrt D,
 \end{gathered}\\
 Z=t+\frac18\log D.
 \label{d:implicit-t}
\end{gather}
The unknowns are $f,Z$.  For any fixed gradient of $f$,
Equation~\eqref{d:implicit-t} determines $t$ uniquely and smoothly:
the derivative of its right side with respect to $t$ is
$1+pv/4>0$, and its range is all of $\R$.

Where $\sigma>0$, set $e=\nabla f/\sigma$ and
$A_j=I-(1-j)e\otimes e$.  Define
\begin{equation}\label{d:ellipticity}
 \begin{gathered}
 \chi=\frac{4d}{4+pv},\qquad
 A_d=I-w\otimes w,\qquad
 A_\chi=I-\frac{p+4}{4+pv}w\otimes w,\\
 H^f=\frac l{\sqrt D}\Hess f,\qquad S_1=K+H^f.
 \end{gathered}
\end{equation}
The latter matrix formulas extend smoothly across $\nabla f=0$.
For a symmetric tensor $A$, $\tr_{A_\chi}A$ means contraction with
$A_\chi$.  The physical equations are
\begin{align}
 \tr_{A_\chi}S_1&=F=h+C,\label{d:trace}\\
 V&=uA_\chi\bigl(4\nabla Z+K(w,\cdot)\bigr),\qquad
 \divg V=u\Xi.\label{d:divergence}
\end{align}
Use $h=b_-$ on $B_b$, $h=b_+$ on $B_w$, and prescribe
\begin{equation}\label{d:boundary}
 \frac{V_\nu}{u}=\mathcal B:=-H+w_\nu(F-P_B)-N_Bd,\qquad
 N_B>1+|H|,
\end{equation}
where $N_B$ is a fixed smooth invariant positive bound on $B$.
On the outer truncation $S_R$, set $f=Z=0$.

Relative to $e$, let $T$ be the transverse two-by-two block of $S_1$,
$M$ its mixed transverse--$e$ block, and $q=S_1(e,e)$.  Let
$T^{\rm tf}=T-\frac12(\tr T)I_\perp$, and use $\nabla_\perp$ for
the transverse gradient.  At zero gradient any axis can be used in
these block formulas; the tensor expressions below show their
independence of that choice.  Write $\dd t,\dd Z$ for differentials,
to distinguish them from the scalar $d=D^{-1}$.

\begin{lemma}[Curvature and boundary identities]\label{d:identities}
For $\bar g=g+l^2\dd f^2$ and $\hat g=e^{4t}\bar g$, and for the
actual trace $F=\tr_{A_\chi}S_1$, one has
\begin{equation}\label{d:curvature}
 \frac12e^{4t}R_{\hat g}
 =8\pi\bigl(\mu+J_0(w)\bigr)+\mathcal T+w(F)-F\tr K
       -u^{-1}\divg V,
\end{equation}
where
\begin{align}\label{d:quadratic}
 \mathcal T={}&\frac12|T^{\rm tf}|^2
 +(M+pa\nabla_\perp t)\cdot(M+(p+2)a\nabla_\perp t)
 +4(p+1)|\dd t|_{A_d}^2\notag\\
 &+(\chi-\chi^2/4)q^2+2(p+1)\chi qa\,\partial_et
       -\frac\chi2Fq+\frac34F^2.
\end{align}
It has the square completion
\begin{align}\label{d:squares}
 \mathcal T={}&\frac12|T^{\rm tf}|^2
  +|M+(p+1)a\nabla_\perp t|^2
  +[4(p+1)-v]|\nabla_\perp t|^2\notag\\
 &+4(p+1)d\left(\partial_et+\frac{\chi aq}{4d}\right)^2
  +\frac34\left(F-\frac{\chi q}{3}\right)^2
  +\frac{4d(9-d)}{3(4+pv)^2}q^2.
\end{align}
In particular,
\begin{equation}\label{d:coercivity}
 |T|^2+|M|^2+dq^2+F^2+|\dd t|_{A_d}^2
       \le\Pi_n\mathcal T.
\end{equation}
Here and below $\Pi_n$ denotes a positive polynomial bound in $n$;
its coefficients may depend on the prepared data, fixed collars and
cutoff profiles, but not on the truncation or the solution.  Different
occurrences may denote different such bounds.  If $\dd t=0$, then,
with $Q_f=|T^{\rm tf}|^2+|M|^2+\chi q^2+F^2$,
\begin{equation}\label{d:floor-norm}
 \tfrac12Q_f\le\mathcal T\le Q_f.
\end{equation}
At any boundary face on which $f$ is constant,
\begin{equation}\label{d:boundary-identity}
 \frac{V_\nu}{u}
 =d(H+4\partial_\nu t)-H+w_\nu(F-P_B),\qquad
 H_{\hat g}=e^{-2t}\sqrt d\,(H+4\partial_\nu t).
\end{equation}
\end{lemma}
\begin{proof}
Consider the auxiliary stationary Lorentz metric
$-l^2(\dd b-\dd f)^2+\bar g$.  Its $b$-slice metric is $g$, its
lapse is $U=l\sqrt D$, its shift is $\beta=Uw$, and its slice tensor is
\[
 k_s=-U^{-1}\operatorname{sym}\nabla\beta
     =-H^f-p(\dd t\otimes w^\flat+w^\flat\otimes\dd t).
\]
Set $\mathfrak B(A,B)=A:B-(\tr A)(\tr B)$,
$P_A=A-(\tr A)g$, and $Q_1=K-k_s$.  The Gauss and normal-expansion
identities, with stationarity, $U\tr k_s=-\divg\beta$, and, for the
future unit normal $n_{\rm L}$,
\[
 \operatorname{Ric}_{\rm Lor}(n_{\rm L},n_{\rm L})
 =-n_{\rm L}(\tr k_s)-|k_s|^2+U^{-1}\Delta U,
\]
give the two expressions for the spacetime scalar density
\begin{equation}\label{d:auxiliary-scalar}
 U\bigl(R_g+\mathfrak B(k_s,k_s)\bigr)
 -2\divg\bigl(\nabla U+(\tr k_s)\beta\bigr)
   =U\bigl(R_{\bar g}-2l^{-1}\bar\Delta l\bigr).
\end{equation}
The right side is the static warped-product expression in the
coordinate $b-f$.  Moreover,
\[
 \bar\Delta l=\frac lU\divg\left(\frac UlA_d\nabla l\right),\qquad
 \nabla U-\frac UlA_d\nabla l=-k_s(\beta,\cdot).
\]
Substitute the constraint equations and
$P_K:\nabla\beta=-U\mathfrak B(K,k_s)$ into
Equation~\eqref{d:auxiliary-scalar}.  The result is
\begin{equation}\label{d:graph-scalar}
 UR_{\bar g}
 =U\left[16\pi(\mu+J_0(w))+\mathfrak B(Q_1,Q_1)\right]
       -2\divg(UP_{Q_1}w).
\end{equation}
The conformal terms in $\frac12e^{4t}R_{\hat g}$ are
$-4\bar\Delta t-4|\dd t|_{A_d}^2$.  Differentiating
Equation~\eqref{d:implicit-t} gives
\begin{gather}\label{d:z-differential}
 4\dd Z=(4+pv)\dd t+aH^f(e,\cdot),\\
 V/u=P_{Q_1}w+4A_d\nabla t+wF,\qquad
 \divg(uw)/u=F+(1-\chi)q-\tr K+2(p+1)w(t).
 \label{d:flux-conversion}
\end{gather}
Combining these formulas leaves the quadratic tensor expression
\begin{align}\label{d:quadratic-tensor}
 \mathcal T={}&\tfrac12\mathfrak B(Q_1,Q_1)
   +2P_{Q_1}(w,\nabla t)+4(p+1)|\dd t|_{A_d}^2\notag\\
 &+F\bigl[F+(1-\chi)q+2(p+1)w(t)\bigr].
\end{align}
This expression is smooth at $w=0$; for example
$(1-\chi)q=\frac{p+4}{4+pv}S_1(w,w)$.
Expanding with $\tr T=F-\chi q$ gives
Equations~\eqref{d:curvature} and~\eqref{d:quadratic}.  Completing the
mixed and normal squares gives Equation~\eqref{d:squares}, hence
Equation~\eqref{d:coercivity}.  When $\dd t=0$, the normal quadratic
form in $(\sqrt\chi q,F)$ has eigenvalues $1$ and $(3-\chi)/4$,
which lie between $1/2$ and $1$, proving
Equation~\eqref{d:floor-norm}.  Finally, at a constant-$f$ face,
$(\Hess f)|_{TB}=(\partial_\nu f)\mathrm{II}_B$.  Substituting
this into the trace equation gives the first boundary identity;
the graph and conformal mean-curvature formulas give the second.
\end{proof}

Choose smooth positive invariant weights $\rho_0,\rho_1$ equal to
$r^{-4},r^{-2\gamma}$ on the far end, where
\begin{equation}\label{d:gamma}
 \frac34<\gamma<1.
\end{equation}
The source is
\begin{equation}\label{d:source}
 \Xi=\delta_0\mathcal T+\rho+\delta_0\tau_0a\sigma
               -m_0(t)\mathcal P,\qquad
 \mathcal P=C_n(\rho_0+v\rho_1).
\end{equation}
Here $0<\delta_0<1/4$ is a sufficiently small fixed number, $C_n$ is
a sufficiently large polynomial in $n$, and $m_0$ is a smooth cutoff
between zero and one, equal to one on $t\le-\eps$ and zero on
$t\ge-\eps+1/n$.  The floor estimate below specifies the order in
which $\delta_0$ and $C_n$ are chosen.

\subsection{The four-stage homotopy}\label{d:homotopy}

The compact-map construction will be performed on finite
truncations.  We specify its equations here so that every a priori
estimate applies to the entire homotopy.  In its first three stages
use $\mathcal T$ with the \emph{actual} trace $F$, and use
\begin{gather}\label{d:homotopy-flux}
 V_\lambda=uA_\chi\bigl(4\nabla Z+\lambda K(w,\cdot)\bigr),\qquad
 \divg V_\lambda=u\Xi,\\
 \frac{(V_\lambda)_\nu}{u}
 =[1-(1-\lambda)j(t)]
 \left[\mathcal B-(1-\lambda)(A_\chi K(w,\cdot))_\nu\right],
 \label{d:homotopy-boundary}
\end{gather}
where $0\le j\le1$ is smooth, zero for $t\le0$, and one for $t\ge1$.
The outer data are always $f=Z=0$.
\begin{enumerate}[label=\arabic*.]
\item Run $\lambda$ from $1$ to $0$, leaving $F=h+C$ and the assigned
inner heights unchanged.  Keep $\lambda=0$ in all later stages.
\item Add gradually to $F=h+C$ a collar term $D_0(x,w,h)$.  It is
smooth, bounded, invariant and nondecreasing in $h$.  On black
collars it is nonpositive and supported where
$w\cdot\nu_s<-1/2$ and $h<b_-+2\eta$; on white collars it is
nonnegative and supported where $w\cdot\nu_s>1/2$ and
$h>b_+-2\eta$, for a fixed sufficiently small $\eta>0$.
Choose its strength so that at the full second-stage endpoint,
at the assigned height on every face and in the limiting directions,
\begin{equation}\label{d:directional-inequalities}
 F(x,\nu)\ge P_B+H,\qquad F(x,-\nu)\le P_B-H
 \qquad(|w|=1,\ \chi=0).
\end{equation}
One of these inequalities is already equality by
Equation~\eqref{d:threshold-boundary}; the other is obtained by the
specified sign of $D_0$ and smooth cutoffs.
\item Extend the assigned heights smoothly to a compactly supported
invariant function $b(x)$, constant on smaller collars.  Run
$\zeta$ from $0$ to $1$, with $h|_B=(1-\zeta)b$ and
\begin{equation}\label{d:third-trace}
 F=h+(1-\zeta)\bigl[C+D_0(x,w,h+\zeta b(x))\bigr]
       +\zeta\bigl[\tr_{A_\chi}K+D_1(x,w)\bigr].
\end{equation}
Near each face take $D_1=A_1w\cdot\nu_s$, with
$A_1>1+\sup_B|H|$, and cut it off in the collars.  At the current
assigned height, $h+\zeta b=b$, so
Equation~\eqref{d:directional-inequalities} persists by convexity in
$\zeta$.  At $\zeta=1$ the trace problem has $f=0$: zero solves it,
and the maximum principle gives uniqueness.
\item Retain that trace problem and scale both the divergence source
and its inner boundary flux to zero by a common factor from $1$ to
$0$.  Since $f=0$, at the final endpoint the second problem has the
unique solution $Z=0$.
\end{enumerate}
In every stage, the derivative of $F$ with respect to $h$, holding
the other variables fixed, is at least one.  This strict height
monotonicity will be used for comparison and for the separate trace
solve.  Every cutoff and auxiliary choice is invariant.

\subsection{Height bounds and exclusion of the floor}\label{d:floor-estimates}

The maximum principle for the trace equation, also without the
divergence equation, gives
\begin{equation}\label{d:height-bound}
 |h|+|F|\le C_0.
\end{equation}
Indeed at a critical point $d=\chi=1$, all gradient-dependent collar
terms vanish, and the trace is strictly increasing in $h$ with
bounded remaining terms.  The constant $C_0$ is independent of $n$,
the truncation and the homotopy parameter.  Far out,
$K=C=D_0=D_1=0$ and $F=h$.  Comparison with
$\pm C(r^{-\gamma}-R^{-\gamma})$ gives
\begin{equation}\label{d:radial-height}
 |h|\le C(r^{-\gamma}-R^{-\gamma})\quad(r\le R).
\end{equation}
At a comparison point with the common radial gradient, the leading
coefficient of the Hessian trace of $r^{-\gamma}$ is
$\gamma[(\gamma+1)\chi-2]<0$, uniformly for $0<\chi\le1$ because
$\gamma<1$; the asymptotic metric errors are smaller.  On the outer
face this implies
$\sigma\le C\tau_0^{-1}R^{-1-\gamma}$.  Since $Z=0$ there,
Equation~\eqref{d:implicit-t} then gives $t>-\eps$ for all sufficiently
large $R$, at fixed $n$.

\begin{lemma}[No contact with the floor]\label{d:floor}
Choose $\delta_0>0$ sufficiently small, then $C_n$ sufficiently large
with polynomial growth.  On every sufficiently large finite
truncation, no smooth solution of any homotopy stage satisfying
$t\ge-\eps$ can attain $t=-\eps$.
\end{lemma}
\begin{proof}
At an inner boundary contact in the first three stages, $j(t)=0$.
Equations~\eqref{d:homotopy-boundary} and~\eqref{d:boundary-identity}
therefore give
$4\partial_\nu t=-H-N_B<0$, incompatible with a minimum when $\nu$
points into the domain.  The preceding outer estimate excludes the
outer face.

Suppose there is an interior contact.  Then $\dd t=0$ and
$\Hess t\ge0$.  The graph Gauss equation in the Riemannian warped
product $g+l^2\dd b^2$ gives
\begin{equation}\label{d:floor-gauss}
 \frac12e^{4t}R_{\hat g}
 \le\frac12R_g-v\Ric_g(e,e)+\mathcal S(H^f),
\end{equation}
where, for a symmetric tensor $A$,
\[
 \mathcal S(A)=\frac14(\tr_\perp A)^2
 -\frac12|A_{\perp\perp}^{\rm tf}|^2
 +dA_{ee}\tr_\perp A-d|A_{e\perp}|^2.
\]
In fact $l^{-1}\Hess l=p\Hess t$ at contact, so the warping term
$-v\tr_\perp(l^{-1}\Hess l)$ and the conformal term
$-4\bar\Delta t$ have favorable signs; $H^f$ is the graph second
form there, up to the immaterial choice of normal.

Directly from Equation~\eqref{d:quadratic}, at $\dd t=0$,
\begin{equation}\label{d:floor-algebra}
 \mathcal T-\mathcal S(S_1)
 =|T^{\rm tf}|^2+(1+d)|M|^2
  +\chi(1+d-\chi/2)q^2-dFq+\frac12F^2.
\end{equation}
Use $\chi\le d$ and
$|dFq|\le\chi q^2/4+(d^2/\chi)F^2$ together with
$d/\chi\le1+n/4$ and Equation~\eqref{d:floor-norm}.  This gives
\[
 \mathcal T-\mathcal S(S_1)
 \ge\tfrac12\mathcal T-(1+n/4)F^2.
\]
Replacing $S_1$ by $H^f=S_1-K$ introduces only products involving
transverse components and $dq$.  Young's inequality and
Equation~\eqref{d:floor-norm} absorb an arbitrarily small fixed multiple
of $\mathcal T$, at polynomial cost in $n$ for $|K|^2$.  Thus, for a
fixed $c_*>0$ independent of $n$,
\begin{equation}\label{d:floor-coercivity}
 \mathcal T-\mathcal S(H^f)
       \ge c_*\mathcal T-\Pi_n(F^2+|K|^2).
\end{equation}
When $K=0$, improve the last error to $\Pi_n vF^2$.  To see this,
if $(1+p)v$ is sufficiently small, $d$ and $\chi$ are uniformly
close to one and the normal quadratic form in
Equation~\eqref{d:floor-algebra} is uniformly positive definite; its
value at $d=\chi=1$ is $\frac32q^2-Fq+\frac12F^2$.
Otherwise $v\ge c/(1+n)$, so the previous error is bounded by a new
polynomial times $vF^2$.

At contact, differentiating $w$ multiplies the vector slot of $H^f$
by $A_d$, and hence
$|\nabla w|\le\Pi_n(|K|+\sqrt{\mathcal T})$.
View $F$ as a smooth function of $(x,w,h,p)$.  Its explicit $x,w$
derivatives have compact support and polynomial bounds on the
range in Equation~\eqref{d:height-bound}; also $\nabla p=0$ at contact.
Since $F_h\ge1$, for any fixed $\eta_0>0$,
\begin{equation}\label{d:floor-height-derivative}
 w(F)\ge\tau_0a\sigma-\eta_0\mathcal T
                  -\Pi_n\mathbf1_{\rm comp}.
\end{equation}
Here $\mathbf1_{\rm comp}$ denotes the indicator of a fixed
sufficiently large compact set; it can equally be replaced by a
fixed compactly supported majorant.
The difference between the physical flux and $V_\lambda$ costs at
most $\eta_0\mathcal T+\Pi_n\mathbf1_{\rm comp}$ after taking its
divergence and dividing by $u$.  Indeed differentiate
$uA_\chi K(w,\cdot)$ using the preceding estimate.  The remaining
factor $\nabla\log u=H^f(w,\cdot)$ contracts against
$A_\chi K(w,\cdot)$ and involves only
$\chi H^f_{ee}$ and $H^f_{e\perp}$, which are controlled by
Equation~\eqref{d:floor-norm}.

Compare Equations~\eqref{d:curvature} and~\eqref{d:floor-gauss}, and use
Equations~\eqref{d:floor-coercivity} and
\eqref{d:floor-height-derivative}.  Outside the fixed compact set,
$K=0$, $F=h=O(r^{-\gamma})$, and $\Ric_g=O(r^{-3})$.
Choose $\eta_0$ and then $\delta_0$ sufficiently small compared with
$c_*$.  The resulting lower bound is
\begin{equation}\label{d:floor-divergence}
 u^{-1}\divg V_\lambda
 \ge2\delta_0\mathcal T+\tau_0a\sigma
          -\Pi_n\bigl[\mathbf1_{\rm comp}+v(\rho_0+\rho_1)\bigr].
\end{equation}
Choose $C_n$ so that $\mathcal P$ dominates the last error plus
$2\rho$ everywhere.  This requires only polynomial growth, since
$\rho_0$ is positive on the compact set.  At the floor $m_0=1$,
and subtracting the prescribed source from this geometric lower
bound gives at least
\[
 \delta_0\mathcal T+(1-\delta_0)\tau_0a\sigma+\rho>0,
\]
contradicting Equation~\eqref{d:homotopy-flux}.

In the fourth stage, $f=0$, $t=Z$, and $v=0$.  As long as the common
scale factor is positive, the inner floor-contact derivative is
that factor times $(-H-N_B)/4<0$.  At an interior floor contact,
$\mathcal T$ is compactly bounded by $K$; enlarge $C_n$ so that
$\delta_0\mathcal T+\rho-\mathcal P<0$ there.  The divergence of
$4u\nabla t$ is nonnegative at a minimum, giving a contradiction.
At scale zero, the unique mixed-problem solution is $Z=0$.  This
also cannot meet the floor.
\end{proof}

\subsection{Separation of heights and the boundary slopes}\label{d:collar-estimates}

All estimates from now on concern solutions with $t\ge-\eps$.
Thus $l\ge\ell$ and $\tau_0/\ell=\ell^{1/2}$ tends to zero
exponentially as $n\to\infty$.

\begin{lemma}[Height separation from the full regions]\label{d:height-separation}
In the first two homotopy stages, for every compact
$C_*\subset\overline\Omega_{\rm red}$ disjoint from the selected
full black region, there are $\eta'>0$ and $n_0$ such that
$h>b_-+\eta'$ on $C_*$ for $n\ge n_0$ and all sufficiently large
outer truncations.  The analogous statement for a compact set
disjoint from the selected full white region is $h<b_+-\eta'$.
These constants are uniform over the first two stages.
\end{lemma}
\begin{proof}
The quantifier \emph{sufficiently large truncations} allows a
threshold $R_0(n)$ depending on $n$.  If the first conclusion fails,
then for every $\eta'>0$ and every $n_0$ there is an $n\ge n_0$ with
failures at arbitrarily large $R$.  Choose successively $\eta'=1/j$,
$n_j\ge\max\{j,n_{j-1}+1\}$ and
$R_j\ge\max\{j,R_{\rm floor}(n_j)\}$ among those failures, with
$R_{\rm floor}(n)$ the outer threshold in Lemma~\ref{d:floor}.
There are points $x_j$ in the fixed compact set with
$h_j(x_j)\le b_-+1/j$.  Pass to $x_j\to x_*$ and take the lower
spatial relaxed limit
$\underline h$, defined by infima over tail indices and shrinking
balls, followed by the lower semicontinuous envelope.  Near a white
face first extend each $h$ continuously by the constant $b_+$ behind
the face.  This does not assert continuity of the relaxed limit.
By construction, $\underline h(x_*)\le b_-$.

Away from black faces, every smooth lower test $\phi$ for
$\underline h$ at a value sufficiently near or below $b_-$, with
$\nabla\phi\ne0$, has oriented level expansion for $K$ at most
$\underline h+C_b$.  Indeed make the touching strict by a
fourth-order perturbation and transfer it to local minima of
$h-\phi$ along a subsequence.  The contacts are interior, because
white boundary values are higher.  There
$\nabla f=\tau_0^{-1}\nabla\phi$, so
\[
 d,\chi\longrightarrow0,\qquad
 \frac{l}{\tau_0\sqrt D}\longrightarrow\frac1{|\nabla\phi|}.
\]
The lower Hessian comparison in the trace equation therefore gives
the level expansion bound.  The second-stage correction has
$D_0\le0$ at these low values, so it has the favorable sign.

We pass from this lower-test statement to exterior supports of
closed sublevels.  Fix $b_-<k<k'$ sufficiently close to $b_-$ and
put $D_k=\{\underline h\le k\}$.  Apply the increasing
reparametrizations
\[
 r_j(s)=\left(1+\exp[-j^2(s-k-j^{-1})]\right)^{-1}.
\]
Their lower relaxed limit on $\underline h$ is the function $u_k$
equal to zero on $D_k$ and one off $D_k$.  At a point outside the
closed set, lower semicontinuity provides a local positive gap from
$k$; on the set itself the unchanged point gives limiting value zero.
Let $\phi$ define an exterior support at $x_0$, with $\phi(x_0)=0$,
$\nabla\phi(x_0)\ne0$, and $D_k$ locally on its nonpositive side.
Choose a small closed ball with $|\phi|<1/4$, and set
$\phi_\eta=\phi-\eta|x-x_0|^4$, where $\eta>0$.
Then $u_k-\phi_\eta$ has its unique minimum zero at $x_0$ and a
positive gap on the ball's boundary.  Minimize the lower
semicontinuous function $r_j(\underline h)-\phi_\eta$ on this ball.
Because $r_j(\underline h(x_0))\le r_j(k)\to0$, the lower relaxed
limit property gives interior minimizers $y_j\to x_0$, with
minimum values $c_j\to0$.  Their transformed contact values
$a_j=\phi_\eta(y_j)+c_j=r_j(\underline h(y_j))$ tend to zero and,
for each finite $j$, lie strictly between zero and one.

On a sufficiently small neighborhood of $y_j$ the smooth function
$\psi_j=r_j^{-1}(\phi_\eta+c_j)$ is therefore a lower test for
$\underline h$.  The inverse is single-valued and strictly
increasing; moreover $\nabla\phi_\eta(y_j)\to\nabla\phi(x_0)\ne0$,
so $\nabla\psi_j(y_j)\ne0$ for large $j$.  Since $r_j(k')\to1$,
one also has $\underline h(y_j)<k'$ for large $j$.
Increasing reparametrization preserves the gradient direction and
the tangential Hessian divided by gradient length, regardless of
the size of its derivatives.  Apply the already proved lower-test
inequality to each fixed $\psi_j$, and then let $j\to\infty$.
The quartic perturbation has zero first and second jets at $x_0$.
Thus every exterior support of $D_k$, away from black faces, has
expansion at most $k'+C_b$.  The original solution limit is taken
first for each fixed test; the reparametrization limit is a separate
subsequent limit.

This property excludes contact with a white face: the reversed
white face itself would be an exterior support, with black
expansion $-c_w>0$.  The radial height bound makes $D_k$ compact on
the end, since $k<0$.  Apply Lemma~\ref{d:weak-support} in the
original connected truncation $\Omega_R$, with the original
boundary as required inner face and $S_R$ as its only designated
outer face.  Its compact set is exactly
$D_k\cup\mathcal B_{c_b}$: neither the white regions nor omitted
bounded components are adjoined.  Every frontier point of this
union outside $\mathcal B_{c_b}$ lies in the reduced exterior
away from white faces, and hence has the just-proved support bound.
At any component of the smooth frontier of $\mathcal B_{c_b}$,
including components facing omitted regions, a support of the
union also supports that black region, and comparison gives the
bound $c_b<k'+C_b$.  Thus the union contains the required inner
collar and meets all hypotheses of the lemma, whether or not it
is connected.  The lemma encloses it in a
smooth trapped region at any threshold slightly larger than
$k'+C_b$ and still negative.  It is therefore contained in the full
black region at that larger threshold.

Right continuity at $c_b=b_-+C_b$ lets us choose a threshold just
above $c_b$ whose full region still misses $C_*$.  Then choose
$k,k'$ close enough to $b_-$ that the preceding enclosure uses a
smaller threshold.  A purported sequence in $C_*$ with limiting
height at most $b_-$ places a limit point in $D_k\cap C_*$, a
contradiction.  This proves uniform strict separation by the usual
subsequence criterion.

For the other color apply the argument to $-h,-K,-C$.  Extend the
original heights constantly by their low boundary values locally
behind black faces before taking the reverse relaxed limit.
The reversed black face is a forbidden support for the high
sublevels.  The ambient domain is now the fixed black complement
in $\Omega_R$, connected after the black pocket filling, with all
black frontiers and $S_R$ designated as outer faces.  Adjoin only
the full white region to the high sublevel.  Supports on any full
white frontier, including hidden components, are controlled by its
smooth threshold equation.  Visible black faces are excluded as
just explained.  A hidden black component has positive distance
from the closure of the reduced domain: a fixed collar along its
connected smooth face lies in a single adjacent complement
component.  Thus the union is disjoint from every designated outer
face.  Omitted components are not adjoined and introduce no other
support frontier.  Any pocket filling in the enclosure lemma is
relative to these designated faces.  Hausdorff right continuity at
$c_w$ gives the same contradiction.  If a selected
color is absent, right continuity means nearby larger-threshold
regions are empty, and the enclosure contradiction is unchanged.
\end{proof}

\begin{proposition}[Collar comparison]\label{d:collars}
For all sufficiently large $n$ and sufficiently large outer
truncations, solutions in the first two homotopy stages satisfy,
with constants $c,C'>0$ independent of $n,R$ and the stage parameter,
\begin{equation}\label{d:collar-slopes}
 \frac c{\tau_0}\le\partial_\nu f\le\frac{C'}{\tau_0}\quad(B_b),\qquad
 \frac c{\tau_0}\le-\partial_\nu f\le\frac{C'}{\tau_0}\quad(B_w).
\end{equation}
They also satisfy $b_-\le h\le b_+$ on any fixed compact core needed
in Equation~\eqref{d:offset-margin}.  In the third stage one has
$|\partial_\nu f|\le C'/\tau_0$.  These conclusions need no upper
bound for $Z$.
\end{proposition}
\begin{proof}
In a black collar, test with $h_0=b_-+\psi(s)$, where
$\psi(0)=0$ and $\psi'>0$.  At a comparison point the gradient is
common to the solution and test, so $t,l,d,\chi$ are computed using
that gradient and the given value of $Z$.  The trace of the test is
\begin{equation}\label{d:collar-trace}
 P_s+aH_s+\chi K(\nu_s,\nu_s)
 +a\chi\left[
  \frac{\Hess s(\nu_s,\nu_s)}{|\nabla s|}
                  +|\nabla s|\frac{\psi''}{\psi'}\right],
\end{equation}
where $H_s$ and $P_s$ use $\nu_s$.  For slopes bounded above and
below by positive constants,
\begin{equation}\label{d:collar-smallness}
 d\le C(\tau_0/\ell)^2=C e^{-n\eps},\qquad
 1-a\le d,\qquad \chi\le d,\qquad
 n\chi=\frac{4nd}{4+pv}\ge2d.
\end{equation}

Use Lemma~\ref{d:height-separation} on the far leaf of a fixed short
collar, and choose a small positive slope so that a lower test stays
below the solution there.  In this direction $D_0$ vanishes.
The stable-leaf inequality and $C=C_b-k_Bs$ give a positive
trace-minus-right-side residual at least $k_Bs/2-O(d)$, apart from
the last logarithmic-slope term in Equation~\eqref{d:collar-trace}.
Choose
\[
 (\log\psi')'=An\quad(0\le s\le1/n),
\]
nonnegative and smoothly cut off to zero by $s=2/n$, with a fixed
sufficiently large $A$.  By Equation~\eqref{d:collar-smallness} its
contribution dominates $O(d)$ on the initial interval.  From
$s\ge1/n$ onward, the $s$ margin dominates the exponentially small
$d$ error.  The slope changes by at most the fixed factor $e^{2A}$;
choose its initial size after $A$, keeping the whole test small.
This constructs a lower barrier with strictly positive boundary
derivative.

All constants in this choice are independent of $n$ and $R$.
The far leaf is fixed with the domain; height separation gives one
fixed $\eta'$ there for every sufficiently large $n$ and
$R\ge R_0(n)$.  Choose $A$ from the fixed geometric residual
constants first, and only then choose the initial slope using
$\eta'$, $k_B$, the collar width and the fixed ratio $e^{2A}$.
Thus these comparisons introduce no arbitrary fixed-$n$ constants
into the polynomial bounds $\Pi_n$.

For an upper barrier choose a large positive slope and the opposite
logarithmic derivative.  The smooth leaf and offset variations are
$O(s)$, whereas the height increase dominates them; the negative
logarithmic-slope term controls the initial $d$ errors.  The global
height bound supplies comparison on the far leaf.  Strict height
monotonicity gives the comparisons throughout the collar.  Equality
at its initial face then gives the two black boundary derivative
bounds in Equation~\eqref{d:collar-slopes}.  Reverse signs for white
faces.  These barriers near the appropriate faces, together with
height separation off their own-color collars, give
$b_-\le h\le b_+$ on the required compact core.

In the third stage, start at the current assigned height
$(1-\zeta)b$ and use lower and upper tests with large negative and
positive slopes, respectively.  Taper the magnitude of each slope
using the negative logarithmic derivative just described.  The
limiting directional inequalities
\eqref{d:directional-inequalities} give the respective residual signs
at the face.  At fixed boundary height the errors are
$O(s+d+\chi)$; in particular, the term $\tr_{A_\chi}K$ is compared
using the common leaf normal.  The height shift controls the $s$
error, and the logarithmic-slope term, using
Equation~\eqref{d:collar-smallness}, controls $d+\chi$.  This gives
both boundary derivative bounds in the third stage.  All uses of
$l\ge\ell$ occur at a common test gradient, so a ceiling for $Z$ is
unnecessary.
\end{proof}
